\documentclass{article}

\usepackage[T1]{fontenc}
\usepackage[utf8]{inputenc}
\usepackage{amsmath}
\usepackage{amssymb}
\usepackage{amsthm}
\usepackage[hidelinks]{hyperref}

\newtheorem{theorem}{Theorem}
\newtheorem{lemma}{Lemma}
\newtheorem{proposition}{Proposition}
\newtheorem{corollary}{Corollary}
\newtheorem{definition}{Definition}

\title{CausalQIF: A Typed Lean Core for Finite Agent-Audit Certificates}
\author{Lucia Yizi Zhang}
\date{}

\begin{document}
\maketitle

\begin{abstract}
CausalQIF is a Lean 4 artifact for finite, typed reasoning about
agent-audit certificates.  The artifact isolates the formal core behind
trace recoverability, conditional data processing, cut-variable certificates,
and d-separation premises over finite probability spaces.  Its purpose is to
make explicit which certificate steps are proved, which are reductions to
finite premises, and which graph-to-statistical bridges remain external
assumptions.  The central contribution is a checked certificate language in
which audit claims can be stated with their premise ledger rather than hidden
inside prose.
\end{abstract}

\section{Introduction}

Agent-audit claims mix three different obligations: information inequalities,
graphical separation assumptions, and deployment-specific measurement claims.
CausalQIF separates these obligations in a typed finite setting.  Entropy,
mutual information, conditional mutual information, and conditional data
processing are expressed over finite PMFs.  Graph-level claims are routed
through explicit d-separation and positive-model interfaces.  Certificate
theorems consume those interfaces as visible premises.

The artifact's punchline is simple: a certificate is only as strong as the
typed assumptions it discharges.  CausalQIF therefore records both successful
mechanization and remaining bridge assumptions.

\section{Formal Scope}

The active Lean package is \texttt{causal\_qif}; the public namespace is
\texttt{CausalQIF}.  The default entry point is:
\[
  \texttt{import CausalQIF}.
\]
The source tree is organized into graph, d-separation, information-theory,
certificate, example, and experimental bridge modules.  The paper-facing source
of truth is the repository root documentation:
\begin{itemize}
  \item \texttt{docs/LEAN.md} for build scope and active assumptions;
  \item \texttt{docs/THEOREM\_DEPENDENCIES.md} for the dependency map;
  \item \texttt{THEOREM\_INDEX.csv} for theorem names, premises, and nonclaims.
\end{itemize}

\section{Certificate Spine}

The formal development supports four reusable certificate components.

\begin{enumerate}
  \item Finite entropy and conditional mutual information primitives.
  \item Conditional DPI: if a probe reaches the action only through the state
  conditional on the trace, then its conditional mutual information with the
  action lower-bounds the true state-action relevance.
  \item Trace-gap decomposition and cut-variable reductions: structural
  upper-budget claims are stated as finite information inequalities once a
  cut/capacity premise is supplied.
  \item Autoregressive zero-cut witnesses: deterministic trace recovery rules
  out positive finite-Shannon residual internal-state witnesses.
\end{enumerate}

\section{Premise Ledger}

The artifact intentionally distinguishes closed finite algebra from external
modeling assumptions.  In particular, conditional DPI is a theorem once the
finite conditional Markov premise is supplied.  Cut-set certificate claims keep
the capacity premise explicit.  General d-separation to conditional
independence is routed through positive finite model assumptions and the
current bridge modules; unresolved graph-to-statistical obligations remain
visible in the documentation rather than being presented as closed theorems.

\section{Artifact Layout}

\begin{verbatim}
CausalQIF/
  lean/
    CausalQIF.lean
    CausalQIF/
      Graph/
      DSeparation/
      InfoTheory/
      Certificates/
      Examples/
      Experimental/
  docs/
  THEOREM_INDEX.csv
  provenance/
\end{verbatim}

\section{Status}

This manuscript is a skeleton for the formal-methods paper.  The active
technical content is the Lean artifact and its theorem ledger.  Deployment
studies are not part of this paper's evidence chain and are not used as support
for the formal claims.
Operator-valued Shannon limits belong to a separate information-theory
manuscript.

\end{document}
